\documentclass[a4paper,10pt]{article}
\usepackage{latexsym}
\usepackage{amsmath}
\usepackage{enumitem}
\usepackage[hidelinks]{hyperref}
\usepackage[top=0.8in,bottom=0.8in,left=0.8in,right=0.8in]{geometry}
\usepackage{titlesec}
\usepackage{parskip}

\titleformat{\section}{\bfseries\uppercase}{\thesection}{1em}{}
\titlespacing*{\section}{0pt}{6pt}{4pt}
\setlength{\parindent}{0pt}
\setlist[itemize]{noitemsep, topsep=2pt, leftmargin=*}

\begin{document}

\begin{center}
    {\LARGE \textbf{Ricardo Digiovanni Frumento}} \\
    \vspace{2mm}
    \href{mailto:ricardofrumento@gmail.com}{ricardofrumento@gmail.com} \hspace{1em} | \hspace{1em}
    +1 813 809 6291 \\
    % \href{https://www.linkedin.com/in/ricardo-digiovanni-frumento-072673108}{LinkedIn} \hspace{1em} | \hspace{1em}
    % \href{https://github.com/RicardoFrumento}{GitHub} \hspace{1em} | \hspace{1em}
    % \href{https://ricardofrumento.github.io}{Website}
\end{center}

\vspace{2mm}

\section*{Professional Summary}
PhD candidate in Computer Science at the University of South Florida specializing in robotic manipulation, multi-object grasping, and learning-based decision making. Researcher in the Robot Perception and Action Lab (RPAL) advised by Dr. Yu Sun, with experience in simulation-driven robotics, benchmarking, control, reinforcement learning, and experimental evaluation. Author of peer-reviewed publications in IEEE RA-L, IEEE T-ASE, and CoRL workshop venues, with additional experience teaching robotics, AI, and deep learning courses.

\section*{Education}
\textbf{University of South Florida} \\
PhD in Computer Science \hfill Aug 2023 -- May 2027 (expected) \\
Relevant Coursework: Deep Reinforcement Learning, Autonomous Robots, Parallel Computing, Deep Learning

BS in Computer Science \hfill Aug 2019 -- May 2023 \\
Relevant Coursework: Artificial Intelligence, Introduction to Robotics, Image Processing, Computational Geometry

\section*{Experience}

\textbf{University of South Florida} \\

\emph{Research Assistant, Robot Perception and Action Lab (RPAL)} \hfill Jan 2023 -- Present
\begin{itemize}
    \item Conduct research in robotic manipulation and multi-object grasping, focusing on grasp planning, scene interaction strategies, and learning-based decision making for cluttered environments.
    \item Design and evaluate simulation and benchmarking workflows for manipulation research, including multi-object grasping studies, scene replication, and comparative evaluation across control and learning methods.
    \item Develop robotics software and experimental pipelines using tools such as Python, MuJoCo, ROS, PyTorch, and OpenCV to support data collection, analysis, and system evaluation.
    \item Contribute to the COMPARE benchmarking ecosystem for reproducible evaluation in robotic manipulation and co-author publications in IEEE RA-L, IEEE T-ASE, and CoRL workshop venues.
    \item Present and discuss research findings in academic settings, collaborating with faculty and students on robotics research, experimentation, and technical documentation.
\end{itemize}

\emph{Teaching Assistant} \hfill Aug 2022 -- Present
\begin{itemize}
    \item Support upper-level computer science courses including Artificial Intelligence, Deep Learning, Robotics, and Automata Theory through office hours, grading, review sessions, and student mentoring.
    \item Mentor 150+ students across multiple semesters, helping reinforce technical concepts in machine learning, robotics, programming, and theoretical computer science.
    \item Develop assessments and instructional support materials to improve student understanding and course delivery.
\end{itemize}

\emph{Peer Leader} \hfill Sep 2021 -- Aug 2022
\begin{itemize}
    \item Facilitated collaborative study groups in introductory computer science courses, contributing to a 12\% improvement in average exam performance.
\end{itemize}

\section*{Selected Technical Projects}

\textbf{Gather-Grasp: Multi-Object Robotic Manipulation Pipeline}
\begin{itemize}
    \item Designed a research pipeline for studying gather-and-grasp strategies in cluttered multi-object manipulation settings.
    \item Evaluated action choices such as pushing, dragging, stacking, and pick-and-place to improve grasp opportunities in simulation.
    \item Built supporting tools for experiment configuration, data collection, and post-hoc analysis of manipulation outcomes.
\end{itemize}

\textbf{COMPARE / Benchmarking for Robotic Manipulation}
\begin{itemize}
    \item Contributed to open-source benchmarking efforts aimed at improving reproducibility and standardization in robotic manipulation research.
    \item Helped define evaluation-oriented workflows for comparing learning and control methods across shared tasks and environments.
\end{itemize}

\textbf{Simulation and Real-System Experimentation}
\begin{itemize}
    \item Developed and maintained robotics experimentation workflows bridging simulation and physical systems, including manipulation scenarios involving robotic arms, grippers, and camera-based perception.
    \item Worked on scene generation, object interaction analysis, and experimental tooling for grasping and manipulation studies.
\end{itemize}

\section*{Publications}
\begin{itemize}
    \item T. Chen, \textbf{R. Frumento}, et al. (2025). \emph{Benchmarking Multi-Object Grasping}. IEEE Robotics and Automation Letters. \href{https://doi.org/10.1109/LRA.2025.3595043}{doi:10.1109/LRA.2025.3595043}
    \item Z. Ye, \textbf{R. Frumento}, Y. Sun. (2025). \emph{Only Pick Once: Efficient Multi-Object Picking}. IEEE Transactions on Automation Science and Engineering. \href{https://ieeexplore.ieee.org/document/10965785}{doi:10.1109/TASE.2025.3560956}
    \item A. Norton, \textbf{R. D. Frumento}, et al. (2024). \emph{Standards for Open-Source Benchmarking in COMPARE}. CoRL 2024 Workshop. \href{https://openreview.net/forum?id=eDRxok2gvr}{OpenReview}
    \item \textbf{R. Frumento}. (2020). \emph{Dynamical Systems and Markov Chains}. \href{https://arxiv.org/abs/2012.12442}{arXiv:2012.12442}
\end{itemize}

\section*{Presentations}
\begin{itemize}
    \item \textbf{Panelist}, 2nd Bellini College REU Symposium -- Spring 2026, University of South Florida, May 1, 2026
    \item \textbf{Presenter}, CoRL 2024 Workshop -- Munich, Germany
    \item \textbf{Attendee}, ICRA 2025 -- Atlanta, USA
\end{itemize}

\section*{Technical Skills}
\textbf{Programming:} Python, C++, C, MATLAB, Fortran \\
\textbf{Robotics / ML:} PyTorch, TensorFlow, OpenCV, ROS, MuJoCo, CoppeliaSim \\
\textbf{Tools:} Git, Docker, Linux/Ubuntu \\
\textbf{Languages:} Portuguese (Native), English (Fluent), Spanish (Fluent), French (Elementary)

\section*{Certifications}
\begin{itemize}
    \item SQL Essential Training -- LinkedIn Learning (2024)
    \item Artificial Intelligence Foundations: Neural Networks -- LinkedIn Learning (2023)
    \item Responsible Conduct of Research -- CITI Program (2022)
\end{itemize}

\end{document}
